
\section{Tautness}
\subsection{Tautness for vf-categories}

Let $\X$ be a vf-category. Note first that if a vf-arrow $u \colon p \vuto\tau q_{\phi t}$ is the thickening of $v \colon p \vuto\sigma q_s$ with respect to some cover $\phi \colon \tau \epi \sigma$, i.e.\ $u = \phi^* v$, then in particular $\psi_1^* u= \psi_2^* v$ for any pair of $\FF$-arrows $\psi_1 ,\psi_2$ such that $\phi \circ \psi_1 = \phi\circ\psi_2$.  

\begin{defn}
    For a cover $\phi \colon \tau \epi \sigma$, a vf-arrow $u \colon p \vuto{\tau} q_{\phi t}$ in $\X$ is \emph{$\phi$-reducible} (or \emph{reducible along $\phi$}) if $\psi_1^* u = \psi_2^*u$ for every pair of parallel $\FF$-arrows $\psi_1,\psi_2$ such that $\phi \circ \psi_1 = \phi\circ\psi_2$. 
\end{defn}
\begin{rmk}
    Equivalently, to see that $u \colon p \vuto{\tau} q_{\phi t}$ is $\phi$-reducible, it suffices to see that $\pi_1^* u = \pi_2^* u$ where $\pi_1,\pi_2 \colon \tau \times_\sigma \tau \rightrightarrows \tau$ is the kernel pair of $\phi$.
\end{rmk}

As we now define, in a \emph{taut} vf-category, reducible vf-arrows can be {uniquely reduced}.

\begin{defn}\label{def:taut}
    A vf-category $\X$ is \emph{taut} if for every $\phi$-reducible vf-arrow $u \colon p \vuto\tau q_{\phi t}$ there exists a unique vf-arrow $u_0 \colon p \vuto\sigma q_s$ such that $u = \phi^* u_0$.
    % \begin{enumerate}
    %     \item for every $\phi$-reducible vf-arrow $u \colon p \vuto\tau q_{\phi t}$ there exists a unique vf-arrow $u_0 \colon p \vuto\sigma q_s$ such that $u = \phi^* u_0$;
    %     % \item the canonical map $\X(p, (q_t)_{t:\coprod_{i} \sigma_i} ) \to \prod_{i} \X(p, (q_s)_{s:\sigma_i})$ defined by specialization along each coproduct injection $\sigma_i \hookrightarrow \coprod_i \sigma_i$ is a bijection. 
    % \end{enumerate} 
\end{defn}

\begin{ex}\label{ex:top-sp-are-taut-vf-posets}
    Topological spaces, seen as vf-posets as in \Cref{ex:top-sp-as-vf-posets}, are taut. 
\end{ex}

\begin{lem}
The vf-category $\Set$ is taut.
\begin{proof}
    Fix a set $X$, an $S$-family of sets $(Y_s)_{s\in S}$, and a filter $\sigma$ on $S$. For a cover $\phi \colon \tau \epi \sigma$, let $u \colon X \to \int_{t:\tau} Y_{\phi t}$ be a function into the reduced product of the family $(Y_{\phi t})_{t\in T}$ which, seen as a vf-arrow, is $\phi$-reducible. We need to show that $u$ factors uniquely through the reindexing function $\phi^* \colon \int_{s:\sigma} Y_s \to \int_{t:\tau} Y_{\phi t}$.
    \[\begin{tikzcd}
        X & {\int_{t:\tau}Y_{\phi t}} \\
        {\int_{s:\sigma}Y_{s}}
        \arrow["u", from=1-1, to=1-2]
        \arrow["{u_0}"', dashed, from=1-1, to=2-1]
        \arrow["{\phi^*}"', from=2-1, to=1-2]
    \end{tikzcd}\]
    Let $\rho \coloneqq \tau\times_\sigma\tau$ be the kernel pair of $\phi$, with its two projections $\pi_1,\pi_2 \colon \rho \rightrightarrows \tau$, and denote by $\gamma \colon \rho \epi \sigma$ the composite $\phi\pi_1=\phi\pi_2$. Since $u$ is $\phi$-reducible, we know that $\pi_1^* u = \pi_2^*u$, so that the claim follows if we show that the following diagram is an equalizer in $\Set$:
\[\begin{tikzcd}
	{\int_{s:\sigma} Y_s} & {\int_{t:\tau}Y_{\phi t}} & {\int_{r:\rho}Y_{\gamma r}}
	\arrow["{\phi^*}", from=1-1, to=1-2]
	\arrow["{\pi_2^*}"', shift right, from=1-2, to=1-3]
	\arrow["{\pi_1^*}", shift left, from=1-2, to=1-3]
\end{tikzcd}\]
To see this, recall first that elements of the reduced product $\int_{s:\sigma}Y_{s}$ are in bijection with $\FF$-arrows $\sigma \to \sum_{s\in S} Y_s$ lying above $S$, as shown below. 
\[\begin{tikzcd}
	&&& {\sum_{s\in S}Y_s} \\
	\rho & \tau & \sigma & S
	\arrow[from=1-4, to=2-4]
	\arrow["{\pi_2}"', shift right, from=2-1, to=2-2]
	\arrow["{\pi_1}", shift left, from=2-1, to=2-2]
	\arrow["\phi", two heads, from=2-2, to=2-3]
	\arrow[dashed, from=2-3, to=1-4]
	\arrow[hook, from=2-3, to=2-4]
\end{tikzcd}\]
The claim follows then recalling that $\phi$ is the coequalizer of its kernel pair since every epimorphism in $\FF$ is effective.
\end{proof}
\end{lem}

\subsection{Vf-fully-faithfulness and subspaces}In this subsection, we introduce some properties of vf-functors which will be central later on, but which are also of independent interest. We will see how they behave as expected under the hypothesis of tautness. Moreover, we will employ them to introduce the notion of \emph{(open) subspaces} of a vf-category.

\begin{defn}
    A vf-functor $F \colon \X \to \mathscr Y$ is \emph{vf-faithful} (resp.\ \emph{vf-fully-faithful}) if, for each point $p$, $S$-family of points $(q_s)_{s\in S}$ and filter $\sigma$ on $S$, the function
    \[ \Hom^{\X}_\sigma (p, q_s) \to \Hom^{\Y}_\sigma (F(p), F(q_s)) \]
    is injective (resp.\ bijective).
\end{defn}

\begin{defn}
    A vf-functor $F \colon \X \to \mathscr Y$ is \emph{vf-full} if for each vf-arrow $u \colon F(x) \vuto{s:\sigma} F(y_s)$ in $\Y$ there exists a cover $\phi \colon \tau \epi \sigma$ in $\FF$ and a vf-arrow $v \colon x \vuto{t:\tau} y_{\phi t}$ in $\X$ such that $F(v) = \phi^*(u)$.
\end{defn}

\gabrielnote{Maybe fist define vf-full, then the factorization vf-full surj-on-ob / ff, then define the topological reflection from the factorization of the morphism into the terminal.}

\begin{lem}
    Let $F \colon \X \to \Y$ be a vf-functor between taut vf-categories. Then, the following are equivalent:
    \begin{enumerate}
        \item $F$ is vf-full and vf-faithful;
        \item $F$ is vf-fully-faithful.
    \end{enumerate}
\end{lem}
\begin{proof}
    Clearly, $(2)\Rightarrow (1)$. Conversely, suppose $F$ is vf-full and vf-faithful. To show that it is vf-fully-faithful we need to show that, for each $p \in \X_0$ and each $(q_s)_{s:\sigma} \in \X_0^\sigma$, the function
    \[ \X( p, (q_s)_{s:\sigma}) \to \Y (F(p), (F(q_s))_{s:\sigma}) \]
    is surjective, so let $u \colon F(p) \vuto\sigma F(q_s)$ be a vf-arrow in $\Y$. By vf-fullness of $F$, there exist a cover $\phi \colon \tau \epi \sigma$ and a vf-arrow $v \colon p \vuto\tau q_{\phi t}$ in $\X$ such that $F(v)= \phi^* u$. Now, for any $\psi_1,\psi_2 \colon \lambda\to \tau$ equalizing $\phi$, we have 
    \[F(\psi_1^*v) = \psi_1^*F(v) = \psi_1^*\phi^*u = \psi_2^*\phi^*u = \psi_2^*F(v) = F(\psi_2^*v)\]
    and hence $\psi_1^*v = \psi_2^*v$ by vf-faithfulness of $F$. In other words, $v$ is $\phi$-reducible: since $\X$ is taut, it follows that there is a (unique) vf-arrow $v_0 \colon p \vuto\sigma q_s$ such that $v = \phi^* v_0$. Thus,
    \[ \phi^* u = F(v) = F(\phi^* v_0) = \phi^* F(v_0)\]
    and hence $u = F(v_0)$ by tautness of $\Y$.
\end{proof}


Recall now that a monomorphism of categories is a functor which is injective on objects and faithful; similarly, a monomorphism of vf-categories is a vf-functor which is injective on points and vf-faithful. Therefore, we could define a \emph{sub-vf-category} of a vf-category $\X$ as a vf-category $\Y$ together with a vf-faithful vf-functor $\Y \to \X$ which is also injective on objects. However, just as how in ordinary category theory we are usually more interested in \emph{full} subcategories, here we also care about a more restrictive notion of sub-vf-category, which we call \emph{subspace}. For instance, tautness will be preserved under subspaces, but not under arbitrary sub-vf-categories.

\begin{defn}
A \emph{subspace} of $\X$ is a vf-category $\Y$ endowed with a vf-fully-faithful functor $\Y \mono \X$ which is also injective on points.
\end{defn}

Concretely, a subspace $\Y$ of $\X$ is determined by a subclass $\Y_0 \subseteq \X_0$ of points, together with the vf-arrows, identities, composites, and reindexings inherited from $\X$. Notationally, we then also write $\Y \subseteq \X$.

\begin{lem}
    Subspaces of taut vf-categories are taut.
    \begin{proof}
        Let $\Y \subseteq \X$ be a subspace of a taut vf-category $\X$. To show that $\Y$ is taut, let $u \colon p \vuto\tau q_{\phi t}$ be a $\phi$-reducible vf-arrow in $\Y$. Since $\Y$ is a subspace of $\X$, $u$ is also a vf-arrow in $\X$, and it is also $\phi$-reducible in $\X$. By tautness of $\X$, there exists a unique vf-arrow $u_0 \colon p \vuto\sigma q_s$ in $\X$ such that $u = \phi^* u_0$. Since both $p$ and $(q_s)_{s\in S}$ lie in $\Y$, it follows that $u_0$ is also a vf-arrow in $\Y$, and hence the claim.
    \end{proof}
\end{lem}


\begin{defn}\label{defn:open-subspace}
A subspace $\Y \subseteq \X$ is \emph{open} if for any point $p$ in $\Y$, any $S$-family of points $(q_s)_{s\in S}$ in $\X$, and any filter $\sigma$ on $S$, the class $\Hom^{\X}_\sigma (p, q_s)$ being inhabited implies that $q_s \in \Y$ for $s:\sigma$. A \emph{neighborhood} of a point $p$ in $\X$ is a subspace of $\X$ containing an open subspace $\Y\subseteq \X$ such that $p \in \Y_0$. 
\end{defn}

\begin{rmk}
Denote by $\bf{\Omega}$ the \emph{vf-category of truth values}, i.e.\ the vf-poset structure defined on the frame $\Omega \coloneqq \mathcal P(1)$ by setting $\alpha \vuleq{\sigma} \beta_s$ if and only if $\alpha \subseteq \int_{s:\sigma} \beta_s$, where 
\[ \int_{s:\sigma} \beta_s \coloneqq \bigcup_{S_0 \in \sigma^\op} \bigcap_{s\in S_0} \beta_s.\]
Then, a subspace $\Y\subseteq \X$ is open if the characteristic map $\X_0 \to {\Omega}$ of $\Y_0 \subseteq \X_0$ can be endowed with the structure of a vf-functor $\X \to \mathbf{\Omega}$.

Assuming classical logic, $\bf{\Omega}$ can be identified with the \emph{Sierpinski vf-category} $\mathbf{2}$. \todo{write this better}
\end{rmk}




\subsection{Tautness for vf-posets}As remarked in \Cref{ex:top-sp-are-taut-vf-posets}, topological spaces are taut as vf-posets. In this subsection, we show that this association yields an equivalence of 2-categories between topological spaces and (small) taut vf-posets, where the inverse construction is determined by the notion of openness of \Cref{defn:open-subspace}. Let us start by showing how to canonically associate a taut vf-poset to each vf-category.

\begin{cons}[The topological reflection]\label{cons:top-refl}
The \emph{topological reflection} $\Tref(\X)$ of a vf-category $\X$ is the vf-poset defined as follows:
\begin{enumerate}
    \item points of $\Tref(\X)$ are the same as those of $\X$;
    \item for each $p \in \X_0$ and $(q_s)_{s:\sigma} \in \X_0^\sigma$, we let $p \vuleq{\sigma} q_s$ hold if and only if there exists a vf-arrow $p \vuto\tau q_{\phi t}$ in $\X$ for some cover $\phi \colon \tau \epi \sigma$ in $\FF$.
\end{enumerate}
% In case $\X$ is itself a vf-poset, we refer to $\eta_{\X}$ as the \emph{tautification map}.
\end{cons}


Both vf-category structures on $\X_0$, namely that of $\X$ and that of $\Tref(\X)$, detect a notion of open subspace: as we now show, in fact, the two coincide.

\begin{lem}
    For a vf-category $\X$, a subclass $\Y_0 \subseteq \X_0$ detects an open subspace of $\X$ if and only if it detects an open subspace of $\Tref(\X)$.
\begin{proof}
    Suppose that $\Y_0$ detects an open subspace $\Y$ of $\X$. To see that $\Y$ is also open as a subspace of $\Tref(\X)$, suppose that $p \vuleq{\sigma} q_s$ in $\Tref(\X)$ with $p \in \Y_0$. By definition, this means that there exists a vf-arrow $p \vuto{\tau} q_{\phi t}$ in $\X$ for some cover $\phi \colon \tau \epi \sigma$; by openness of $\Y$, this entails that $q_{\phi t} \in \Y_0$ for $t :\tau$. More explicitly:
    \[ \set{ t | q_{\phi t} \in \Y_0 } \in \tau \]
    from which, since $\phi$ is a cover in $\FF$:
    \[ \phi \left( \set{ t | q_{\phi t} \in \Y_0 }\right) = \set{ \phi t | q_{\phi t} \in \Y_0 } \in \sigma \]
    and hence $\set{ s | q_s \in \Y_0} \in \sigma$, meaning that $q_s \in \mathscr{Y}_0$ for $s: \sigma$. The converse implication is immediate.
\end{proof}
\end{lem}

In other words, the 2-functor $\Open \colon \vfCat \to \TopSp$ --- mapping each small vf-category to the topological space determined by its open subspaces --- factors through vf-posets via the topological reflection $\Tref$. It is then straightforward to see that this yields two (reflective) 2-adjunctions:
\[\begin{tikzcd}
	\TopSp && \vfPos && \vfCat
	\arrow[""{name=0, anchor=center, inner sep=0}, "\Tinj"', curve={height=12pt}, hook, from=1-1, to=1-3]
	\arrow[""{name=1, anchor=center, inner sep=0}, "{\Open}"', curve={height=12pt}, from=1-3, to=1-1]
	\arrow[""{name=2, anchor=center, inner sep=0}, curve={height=12pt}, hook, from=1-3, to=1-5]
	\arrow[""{name=3, anchor=center, inner sep=0}, "\Tref"', curve={height=12pt}, from=1-5, to=1-3]
	\arrow["\dashv"{anchor=center, rotate=-90}, draw=none, from=1, to=0]
	\arrow["\dashv"{anchor=center, rotate=-90}, draw=none, from=3, to=2]
\end{tikzcd}\]

As we now show, the 2-adjunction on the left corestricts to an equivalence between topological spaces and \emph{taut} vf-posets. Recall indeed by \Cref{ex:top-sp-are-taut-vf-posets} that, for each topological space $X$, the vf-poset $\Tinj(X)$ is taut: thus, to conclude that the adjunction corestricts to an equivalence on taut vf-posets, it suffices to show that $\Tinj (\Open Y) \cong Y$ for every taut vf-poset $Y$. The proof of this fact, \Cref{prop:taut-vf-posets-come-from-spaces}, will make use of the following result: intuitively, taut vf-posets admit `canonical' vf-arrows, corresponding to \emph{neighborhood filters}, which are \emph{minimal} with respect to reindexing.

\begin{prop}\label{prop:canonical-convergence-neighborhoods}
    Let $Y$ be a taut vf-poset and let $p\in Y_0$. Then, there exists a filter $\nu_p$ on $Y_0$ such that:
    \begin{enumerate}
        \item $p \vuleq{\nu_p} y$, where the family $y \colon \nu_p \hookrightarrow Y_0$ is represented by the identity function;
        \item $p \vuleq{\sigma} q_s$ holds in $Y$ if and only if $q \colon \sigma \to Y_0$ defines an $\FF$-arrow $\sigma \to \nu_p$.%  satisfying $(q_s)_{s:\sigma} = \phi^*(y)_{y:\nu_p}$.
    \end{enumerate}  
    Moreover, $\nu_p$ can be described as the filter of {neighborhoods} of $p$ in $Y$.
\begin{proof}
    First, we construct a filter $\nu_p$ with the desired properties. Let $N$ be the set of filters $\nu$ on $Y_0$ such that $p \vuleq{y:\nu} y$: note that $N$ is inhabited at least by the principal filter $\delta_p$ at $p\in Y_0$, as witnessed by the diagonal vf-arrow $\diag^{\delta_p}_p$.  Every filter $\nu \in N$ comes equipped with a canonical map $\nu \hookrightarrow Y_0$, represented by the identity function; thus, we also get a canonical map $\coprod_{\nu \in N} \nu\to Y_0$. Let $\nu_p$ be the filter on $Y_0$ arising via its image factorization.
    \[\begin{tikzcd}[sep = small]
    	{\coprod_{\nu \in N} \nu } && {Y_0} \\
    	& {\nu_p}
    	\arrow[from=1-1, to=1-3]
    	\arrow[two heads, from=1-1, to=2-2]
    	\arrow["y"', hook, from=2-2, to=1-3]
    \end{tikzcd}\]
    That $p \vuleq{y:\nu_p} y$ follows then by tautness, since $p \vuleq{y:\nu} y$ for each $\nu \in N$ implies that $p \vuleq{(\nu,y) : \coprod_\nu \nu} y$ by merging, and hence $p \vuleq{y:\nu_p} y$ as reducibility trivializes in vf-posets.

    To show property (ii), suppose that $p \vuleq{\sigma} q_s$ for some $q \colon \sigma \to Y_0$ and let $q_*(\sigma)$ be the filter in its image factorization. By tautness, $p \vuleq{\sigma} q_s$ entails that $p \vuleq{q_*(\sigma)} y$, i.e.\ that $q_*(\sigma) \in N$. Therefore, considering the inclusion $q_*(\sigma) \mono \coprod_{\nu\in N} \nu$, into the coproduct, we have that the following diagram commutes:
\[\begin{tikzcd}
	\sigma &&& \\
	{q_*(\sigma)} & {\coprod_{\nu \in N}\nu } && {Y_0} \\
	&& {\nu_p}
	\arrow[two heads, from=1-1, to=2-1]
	\arrow["q", curve={height=-18pt}, from=1-1, to=2-4]
	\arrow[from=2-1, to=2-2, hook]
	\arrow[bend left=20, hook, from=2-1, to=2-4]
	\arrow[curve={height=12pt}, from=2-1, to=3-3]
	\arrow[from=2-2, to=2-4]
	\arrow[two heads, from=2-2, to=3-3]
	\arrow["y"', hook, from=3-3, to=2-4]
\end{tikzcd}\]
In particular, the composite $\FF$-arrow $\sigma \to \nu_p$ is represented by $q$. % above satisfies $(q_s)_{s:\sigma} = \phi^*(y)_{y:\nu_p}$.

Finally, we describe $\nu_p$ explicitly as the filter of neighborhoods of $p$.
    \begin{itemize}
        \item Let $U \subseteq Y$ be an open neighborhood of $p$. Then, since $p \vuleq{\nu_p} y$ by property (i), by openness it follows that $y \in U_0$ for $y:\nu_p$, i.e.\ $U_0 \in \nu_p$.
        \item Conversely, let $V \subseteq Y$ be a subspace such that $V_0 \in \nu_p$, which equivalently means that $y \in V_0$ for $y: \nu_p$. Note that, by property (ii), $p \vuleq{\delta_p} p$ entails that the identity function realizes an $\FF$-arrow $\delta_p \to \nu_p$; thus, it induces a function $V^{\nu_p} \to V_0^{\delta_p}$. Since $(y)_{y:\nu_p} \in V_0^{\nu_p}$, it follows that $(y)_{y:\delta_p} \in V_0^{\delta_p}$, which is equivalent to $p \in V_0$.
        
        Consider then the subspace $U \subseteq V$ defined by:
        \[ U_0 \coloneqq \set{ q \in V_0 | V_0 \in \nu_q } \]
        By the above, $p \in U_0$; we now show that $U$ is open in $Y$, which proves that $V$ is a neighborhood of $p$. Suppose that $q \vuleq{\sigma} r_s$ in $Y$ with $q \in U_0$: we need to show that $r_s \in U_0$ for $s:\sigma$, which explicitly means that $\set{ s | r_s \in U_0}\in \sigma$. Taking composites, we have 
        \[ q \vuleq{\sigma} (r_s \vuleq{\nu_{r_s}} y ), \mbox{ that is, }q \vuleq{\sum_{s:\sigma}\nu_{r_s}} y,\]
        where now $y$ also denotes the projection $\sum_{s:\sigma}\nu_{r_s} \to Y_0$.  By property (ii) we then know that $y$ also defines an $\FF$-arrow $\sum_{s:\sigma}\nu_{r_s} \to \nu_q$; since $q \in U_0$, which by construction means that $V_0\in \nu_q$ i.e.\ $(y)_{y:\nu_q} \in V^{\nu_q}_0$, applying the induced function $V_0^{\nu_q} \to V_0^{\sum_{s:\sigma} \nu_{r_s}}$ we conclude that 
        \[\set{ s | (y)_{y:\nu_{r_s}} \in V_0^{\nu_{r_s}} }  = \set{ s | V_0 \in \nu_{r_s}} = \set{ s | r_s \in U_0}\in \sigma.\]
    \end{itemize}  
\end{proof}
\end{prop}


\begin{prop}\label{prop:taut-vf-posets-come-from-spaces}
    Let $Y$ be a taut vf-poset. For every point $p$, $S$-family of points $(q_s)_{s\in S}$, and filter $\sigma$ on $S$, the following are equivalent:
    \begin{enumerate}
        \item $p \vuleq{\sigma} q_s$ in $Y$;
        \item $p \vuleq{\sigma} q_s$ in $\Tinj( \Open  Y)$, i.e.\ for every open neighborhood $U \subseteq Y$ of $p$ it holds that $q_s \in U_0$ for $s:\sigma$.
    \end{enumerate}
\begin{proof}
    \looseness=-1
    Clearly $(i)\Rightarrow(ii)$ by definition of openness, so suppose $(ii)$. By \Cref{prop:canonical-convergence-neighborhoods} we know that $q$ realizes an $\FF$-arrow $\sigma \to \nu_p$ where $\nu_p$ is the filter of neighborhoods of $p$ in $\Tinj(\Open Y)$. Let $\nu_p'$ be the filter of neighborhoods of $p$ in $Y$: if we show that the identity map realizes an $\FF$-arrow $\nu_p \to \nu_p'$, i.e.\ that $\nu_p' \subseteq \nu_p$, then by \Cref{prop:canonical-convergence-neighborhoods} we can conclude that $p \vuleq{\sigma} q_s$ in $Y$. In turn, $\nu_p' \subseteq \nu_p$ follows since  open subspaces of $Y$ are open in $\Open Y$ by definition, and the opens of any topological space $X$ are open as subspaces of $\Tinj (X)$.
\end{proof}
\end{prop}

\begin{cor}
    $\TopSp \simeq \vfPostaut$.
\end{cor}

Crucially, topological reflections of small vf-categories are taut \UT{check}: combining this with the previous result, we thus obtain the following characterization.

\begin{cor}\label{prop:convergence-in-topological-reflection-is-witnessed}
    Let $\X$ be a small vf-category. For every point $p$, $S$-family of points $(q_s)_{s\in S}$, and filter $\sigma$ on $S$, the following are equivalent:
    \begin{enumerate}
        \item $p \vuleq{\sigma} q_s$ in $\Tref(\X)$;
        \item for every open neighborhood $\Y \subseteq \X$ of $p$ it holds that $q_s \in \Y_0$ for $s:\sigma$.
    \end{enumerate}
\end{cor}


% \begin{cor}
% Let $\X$ be a small vf-category and let $p \in \X_0$. Let $\nu_p$ be the filter on $\X_0$ defined by the open neighborhoods of $p$. Then, $p \vuleq{x:\nu_p}x$ in $\Tref(\X)$. 
% \begin{proof}
%     By \Cref{prop:convergence-in-topological-reflection-is-witnessed}, it suffices to see that every open neighborhood $\Y$ of $p$ in $\X$ satisfies $(x)_{x:\nu_p} \in \Y_0^{\nu_p}$, which is trivial since $\Y_0$ lies in $\nu_p$.
% \end{proof}
% \end{cor}




